\documentclass[10pt,a4paper]{amsart}
\usepackage[margin=19mm]{geometry}
\usepackage{amsmath,amssymb,mathtools}
\usepackage[scr=rsfs,cal=euler]{mathalfa}
\usepackage{xfrac}
\usepackage[unicode,hidelinks,bookmarks=false]{hyperref}
\newcommand{\ie}{i.e.\ }
\newcommand{\cf}{cf.\ }
\newcommand{\resp}{respectively\ }
\newcommand{\Subsection}{\subsection}
\providecommand{\nobreakdash}{\nobreak\hbox{-}}
\setlength{\emergencystretch}{2em}
\multlinegap0pt
\newtheorem{theorem}{Theorem}
\title{Dynamical analysis of r-Chialvo neuron map with cosine memristive}
\author{Ajay Kumar, V.V.M.S. Chandramouli}
\date{Source: arXiv:2605.00442v1 (CC BY 4.0)}
\begin{document}
\maketitle

\noindent\emph{Source and licence.} Dynamical analysis of r-Chialvo neuron map with cosine memristive. Version arXiv:2605.00442v1 (CC BY 4.0), distributed by arXiv under the Creative Commons Attribution 4.0 International licence, http://creativecommons.org/licenses/by/4.0/, as recorded in the arXiv licence metadata for this exact version and shown on the abstract page https://arxiv.org/abs/2605.00442v1. Full licence notice: \texttt{licenses/arxiv-2605.00442-CC-BY-4.0.txt}.\par\smallskip

\noindent\textit{Editorial scope.} Counted unit: the article's theorem theorem (source0.tex lines 324--326). The article's complete original proof of it is delivered with it (source0.tex lines 329--374, one \texttt{proof} environment of 46 lines). No mathematics is written, repaired or re-derived: the statement and the proof are the article's own lines, in the article's own notation, and no retained line is substituted or re-flowed.  No further source-local context is needed: every cross-reference of every retained line resolves inside the two delivered spans.  Nothing was omitted from any retained span, so the package carries no omission record.  No retained line contains a citation, so the wrapper carries no bibliography and no bibliography entry.  No retained line contains a figure environment, an image-inclusion command or drawing code, so no image is delivered and the package declares no asset dependency.  Editorial formatting only, in addition: an AMS \texttt{amsart} preamble that restates the article's own declarations and macros used by the retained lines, namely the environments \texttt{theorem} on separate counters, together with the standard packages those lines need.  The retained environments print in this document as Theorem 1 in order of appearance, whereas the article prints the same items with the numbering of its own full document.  The complete native source of the article as distributed in the arXiv e-print for this exact version is bundled unchanged as \texttt{sources/199562/source0.tex}.
\begin{theorem}
Let $\mathcal{M}_{r,k}(x,\phi)$ be the map, satisfying the condition $0<k_0<N$ for $N\in \mathbb{N}$, then the map exhibit a unique positive fixed point $(x^*,\phi^*)\in [0,M_1]*[0,M_2]$, whenever the parameter satisfies the following conditions: $k\in (-\frac{1}{M_1},\frac{1}{M_1})$, $0<k_1<1$, $k_2\geq4\pi$ and $r<0$.
\end{theorem}

\begin{proof}
     The fixed points of the map $\mathcal{M}_{r,k}(x,\phi)$ are obtained by solving the following equations simultaneously.
\begin{equation}\label{5a}
        x^2e^{(r-x)}+k_0+kx \cos(\pi\phi)=x,
\end{equation}
\noindent and
  \begin{equation}\label{5b}
         k_1x-k_2\phi=\phi.
  \end{equation}

\noindent From Eq.\,\eqref{5b}, we write
\begin{equation}\label{6}
    \phi=\frac{k_1x}{1+k_2}.
\end{equation}

\noindent Use the value of $\phi$ in Eq.\eqref{5a}, we obtain
\begin{equation}\label{7}
    x^2e^{(r-x)}+k_
0+kx \cos \Big\{\pi\left(\frac{k_1x}{1+k_2}\right) \Big\}=x.
\end{equation}

\noindent Let us rewrite the above expression in the form of a function
\begin{equation}\label{8}
   F(x)= f(x)-x, \hspace{0.4cm}  where \hspace{0.4cm} f(x)=x^2e^{(r-x)}+k_
0+kx \cos \Big\{\pi\left(\frac{k_1x}{1+k_2}\right) \Big\}.
\end{equation}

\noindent Note that the $x$-component of fixed points of the map $\mathcal{M}_{r,k}(x,\phi)$ is same as the roots of above function $F(x)$. Also, the $\phi$-component can be determined by using Eq. \eqref{5b}.

\noindent Assume that  $0<k_0<N$, then it follows that \; $F(0)=k_0 \implies F(0)>0$.


\noindent Further assume that the condition  $r<0$, $M_1=N+2$, $k\in (-\frac{1}{M_1},\frac{1}{M_1})$, $0<k_1<1$, $k_2\geq2\pi$ holds, then we obtain that \; $F(M_1)<0$.

\noindent Therefore, $F(x)$ has at least one positive real root in the interval $[0,M_1]$. Further, to establish the uniqueness of the positive fixed point, we calculate

$F'(x)=(2-x)xe^{r-x}+k \cos\left(\frac{\pi{k_1}x}{1+k_2}\right)-\frac{\pi{k}xk_1}{1+k_2} \sin\left(\frac{k_1{\pi}x}{1+k_2}\right)-1$

\noindent Since $(2-x)xe^{r-x}<\frac{1}{2}$, $-\frac{1}{M_1}<k \cos\left(\frac{\pi{k_1}x}{1+k_2}\right)<\frac{1}{M_1}$,\\ and \;
$\frac{1}{4}<-\frac{\pi{k}xk_1}{1+k_2} \sin\left(\frac{k_1{\pi}x}{1+k_2}\right)<\frac{1}{4}$. 

\vspace{0.2cm}
\noindent So, we get $F'(x)<0$,  for all $x\in[0,M_1]$.

\noindent Hence the map $\mathcal{M}_{r,k}(x,\phi)$ has a unique positive fixed point in $[0,M_1]*[0,M_2]$.
\end{proof}

\end{document}
